\documentclass[a4paper,11pt]{ltjsarticle}
\usepackage{amsmath,geometry,makeidx,etoolbox}
\geometry{left=25mm,right=25mm,top=30mm,bottom=30mm}
\makeindex
\usepackage[colorlinks=true,linkcolor=blue,hypertexnames=false]{hyperref}
\newcommand{\sectionlinkstar}[1]{\section*{\hyperlink{toc}{#1}}\addcontentsline{toc}{section}{#1}}
\renewcommand{\indexname}{索引}
\makeatletter
\patchcmd{\theindex}{\twocolumn[\section*{\indexname}]}{\twocolumn[\sectionlinkstar{\indexname}]}{}{\PackageError{prototype}{Index heading patch failed}{}}
\makeatother
\begin{document}
\hypertarget{toc}{}\tableofcontents
\section*{Test page 1}
物理PDFページと表示ページの両方が1である。
\index{alpha@Alpha（アルファ）}
\index{range@Range（範囲）|(}
\index{Kahler@K\"ahler manifold（K\"ahler多様体）}
\newpage
\section*{Test page 2}
物理PDFページと表示ページの両方が2である。
\index{alpha@Alpha（アルファ）}
\index{beta@Beta（ベータ）}
\index{BRST@BRST cohomology（BRSTコホモロジー）}
\index{BRST@BRST cohomology（BRSTコホモロジー）!physical@physical states（物理状態）}
\begin{equation}a=b\tag{1.1}\label{eq:first}\end{equation}
\newpage
\section*{Test page 3}
物理PDFページと表示ページの両方が3である。
\index{range@Range（範囲）|)}
\index{zeta@Zeta（ゼータ）}
\index{Wess Zumino Novikov Witten@Wess--Zumino--Novikov--Witten model（Wess--Zumino--Novikov--Witten模型）}
\begin{equation}c=d\tag{1.1}\label{eq:second}\end{equation}
\clearpage\printindex
\end{document}
